\subsection{隣接4項間漸化式型}\label{節：隣接4項間漸化式型}
隣接2項間や隣接3項間の漸化式を今まで扱った.
これをより広い範囲に適用してみよう.
例えば
\[
  a_{n+3}+pa_{n+2}+qa_{n+1}+ra_n=0
\]
のような漸化式を考える.
これは三階の同次線形漸化式なので,\sectionref*{節：一次独立・基底・次元}で述べたように,
初めの三項を自由に選ぶ解全体は3次元である.
ここでは一つの補助数列を先に求めることで,残りの計算を二階の式へ分ける.
\sectionref*{節：補助数列の設計}で述べたように,先に単純な更新を決め,それをする量を探そう.
3項をまとめた$d_n=a_{n+2}+\alpha a_{n+1}+\beta a_n$を作り,
$d_{n+1}=\gamma d_n$となることを目指す.
その条件は
\BookMathLayout{\[
  a_{n+3}+\alpha a_{n+2}+\beta a_{n+1}=\gamma \left(a_{n+2}+\alpha a_{n+1}+\beta a_{n}\right)
\]}{
\[
\begin{aligned}
a_{n+3}+\alpha a_{n+2}+\beta a_{n+1}\\
{}=\gamma\left(a_{n+2}+\alpha a_{n+1}+\beta a_n\right)
\end{aligned}
\]
}
である.両辺を左側に集め,元の漸化式と係数を比べると,
\[
  \alpha-\gamma=p,\qquad
  \beta-\gamma\alpha=q,\qquad
  -\gamma\beta=r
\]
となる.最初の二式から
\[
  \alpha=p+\gamma,\qquad
  \beta=q+p\gamma+\gamma^2
\]
なので,最後の式へ代入すると
\[
  \gamma^3+p\gamma^2+q\gamma+r=0
\]
を得る.この三次方程式の実数解を一つ選べば,$\alpha,\beta$も決まる.
三次方程式は少なくとも一つ実数解を持つが,計算しやすい値を見つけられるかは別の問題である.
以下の例題では因数分解や階差の形を手掛かりに選ぶ.
これによって,\,$d_n=a_{n+2}+\alpha a_{n+1}+\beta a_n$とおくと$d_{n+1}=\gamma d_n$が成り立つので,
\[
  a_{n+2}+\alpha a_{n+1}+\beta a_n=d_1\gamma^{\,n-1}
\]
という3項間漸化式が得られる.
ただし,\,\sectionref{節：等差二次特性方程式型}で扱ったのは右辺が定数$r$の場合であり,ここでは右辺が指数関数$d_1\gamma^{n-1}$になる点が異なる.
この形は\sectionref{節：階差等比複合型}の考え方を3項間に応用することで解くことができる.
ここでさらに同じ更新をする既知の量を引くと,右辺の指数関数を消せる場合がある.
例えば第1問では,$p=-3,q=3,r=-1$なので上の方程式は$(\gamma-1)^3=0$となり,
$\gamma=1,\alpha=-2,\beta=1$を得る.
補助数列が$d_n=a_{n+2}-2a_{n+1}+a_n$となる理由は,この係数の一致にある.
元の式が三階差の形であると気付けば,同じ量を階差からも作れる.

  次の数列の一般項を求めてみよう.

  \noindent \bookblocklink{example-four-term-1}{problem}{answer}{【例題】}
  \begin{enumerate}[label=(\arabic*),parsep=3pt,format=\bookitemlink{example-four-term-1}{problem}{answer}]
    \item $\displaystyle a_1=1,\,a_2=2,\,a_3=4$\\
          $a_{n+3}=3a_{n+2}-3a_{n+1}+a_n\quad(n\ge1)$
    \item $\displaystyle a_1=2,\,a_2=1,\,a_3=5$\\
          $a_{n+3}=2a_{n+2}+a_{n+1}-2a_n\quad(n\ge1)$
    \item $\displaystyle a_1=1,\,a_2=4,\,a_3=11$\\
          $a_{n+3}=5a_{n+2}-7a_{n+1}+3a_n\quad(n\ge1)$
  \end{enumerate}

  \noindent\bookblocklink{example-four-term-1}{hint}{problem}{【方針】}
  \begin{enumerate}[label=(\arabic*),parsep=3pt,format=\bookitemlink{example-four-term-1}{hint}{problem}]
    \item $d_n=a_{n+2}-2a_{n+1}+a_n$とおき,階差をさらにとると一定になることを利用します.
    \item $d_n=a_{n+2}-a_{n+1}-2a_n$とおくと$d_{n+1}=d_n$となります.初期条件から$d_n=0$を導き,3項間漸化式として解きます.
    \item $d_n=a_{n+2}-2a_{n+1}+a_n$とおくと$d_{n+1}=3d_n$となります.得られた右辺を満たす数列を引き,階差が一定の形に帰着させます.
  \end{enumerate}

  \noindent\bookblocklink{example-four-term-1}{answer}{problem}{【解答】}
  \begin{enumerate}[label=(\arabic*),parsep=3pt,format=\bookitemlink{example-four-term-1}{answer}{problem}]
    \item $\displaystyle a_n=\frac{n^2-n+2}{2}$
    \item $\displaystyle a_n=2^{n-1}+(-1)^{n-1}$
    \item $\displaystyle a_n=n-1+3^{n-1}$
  \end{enumerate}

  \noindent\bookblocklink{example-four-term-1}{solution}{problem}{【解法】}
  \begin{enumerate}[label=(\arabic*),parsep=3pt,format=\bookitemlink{example-four-term-1}{solution}{problem}]
    \item 漸化式を移項すると,
          \[
            a_{n+3}-3a_{n+2}+3a_{n+1}-a_n=0
          \]
          である.ここで
          \[
            d_n=a_{n+2}-2a_{n+1}+a_n
          \]
          とおくと,
          \[
            d_{n+1}-d_n=a_{n+3}-3a_{n+2}+3a_{n+1}-a_n=0
          \]
          より$d_n$は一定である.初期条件から$d_1=4-2\cdot2+1=1$なので,
          \[
            a_{n+2}-2a_{n+1}+a_n=1
          \]
          となる.さらに$e_n=a_{n+1}-a_n$とおくと$e_{n+1}-e_n=1$であり,$e_1=1$だから$e_n=n$である.したがって,
          \[
            a_n=a_1+\sum_{k=1}^{n-1}e_k
                =1+\sum_{k=1}^{n-1}k
                =\frac{n^2-n+2}{2}
          \]
          を得る.
    \item
          \[
            d_n=a_{n+2}-a_{n+1}-2a_n
          \]
          とおくと,
          \[
            d_{n+1}-d_n=a_{n+3}-2a_{n+2}-a_{n+1}+2a_n=0
          \]
          である.また,$d_1=5-1-2\cdot2=0$だから$d_n=0$であり,
          \[
            a_{n+2}-a_{n+1}-2a_n=0
          \]
          となる.特性方程式$x^2-x-2=0$の解は$2,-1$なので,
          \[
            a_n=A\,2^{n-1}+B(-1)^{n-1}
          \]
          とおける.初期条件より$A+B=2,\,2A-B=1$であるから$A=B=1$を得る.
    \item
          \[
            d_n=a_{n+2}-2a_{n+1}+a_n
          \]
          とおくと,与えられた漸化式から$d_{n+1}=3d_n$である.初期条件より$d_1=11-2\cdot4+1=4$なので,
          \[
            a_{n+2}-2a_{n+1}+a_n=4\cdot3^{n-1}
          \]
          となる.一方,$3^{n-1}$の2階差も$4\cdot3^{n-1}$である.そこで$b_n=a_n-3^{n-1}$とおくと,
          \[
            b_{n+2}-2b_{n+1}+b_n=0
          \]
          となる.$b_1=0,\,b_2=1$より$b_n=n-1$であるから,
          \[
            a_n=b_n+3^{n-1}=n-1+3^{n-1}
          \]
          を得る.
  \end{enumerate}
